% =============================================================
% GECCO 2026 — Theoretical Section Draft (Lyra's half)
% For Claudius to review and integrate into the GECCO branch.
% NOT a standalone document — intended to replace/augment
% Section 2 and relevant parts of Sections 4.2 and 6.2.
% =============================================================

% --- New content begins here ---
% Insert after current Section 2.3 (Composition Levels),
% before Section 3 (Experimental Setup).

\subsection{The Island Functor and Migration Topology}
\label{sec:island-functor}

The island model lifts a single-population strategy to a
multi-population system.  We formalize this lifting as a functor
parameterized by a migration graph.

Let $G = (V, E)$ be an undirected graph on $n = |V|$ vertices.  The
\emph{island functor}
$\mathcal{I}_G : \mathbf{Kl}(M) \to \mathbf{Kl}(M^n)$
acts on objects by
$\mathcal{I}_G(\texttt{Pop}\;a) = (\texttt{Pop}\;a)^n$
and on morphisms by applying the single-island pipeline independently
to each island, then redistributing individuals along edges of~$G$.

Whether $\mathcal{I}_G$ preserves Kleisli composition depends on~$G$:
\begin{itemize}[nosep]
  \item \emph{No edges} ($G$ discrete): migration is the identity and
    $\mathcal{I}_G(\sigma_1 \mathbin{>\!\!=\!\!>} \sigma_2) =
     \mathcal{I}_G(\sigma_1) \mathbin{>\!\!=\!\!>} \mathcal{I}_G(\sigma_2)$
    exactly---the functor is \textbf{strict}.
  \item \emph{Any edges}: migration after each composed step differs
    from migration after their composition.  The discrepancy is the
    \emph{laxator}
    $\varphi_G : \mathcal{I}_G(\sigma_1) \mathbin{>\!\!=\!\!>}
     \mathcal{I}_G(\sigma_2) \Rightarrow
     \mathcal{I}_G(\sigma_1 \mathbin{>\!\!=\!\!>} \sigma_2)$,
    a natural transformation measuring how much interleaved migration
    changes the outcome relative to end-of-run migration.
\end{itemize}

\noindent The laxator controls the \emph{coupling effect} of migration.
It does not change \emph{whether} populations converge---on a unimodal
landscape like OneMax, all topologies reach the optimum.  It controls
\emph{how} they converge: the shape of the transient trajectory, the
rate of diversity loss, and the distribution of genetic material across
islands during the critical early generations.  This is exactly the
diversity-fingerprint phenomenon of Section~4.1, lifted from individual
strategies to multi-population systems.

\begin{remark}[Status of the laxator]
The explicit construction of $\varphi_G$---giving its components and
verifying coherence conditions---remains the central open problem of
this framework.  We take a programmatic approach: identify the laxator
as the key quantity, predict its behavior via graph-theoretic
invariants, and confirm empirically.  The Strict/Lax Dichotomy
(Observation~\ref{obs:dichotomy}) is a first structural theorem about
this quantity: it is the identity if and only if no migration occurs.
\end{remark}


\subsection{From Algebraic Connectivity to Cycle Rank}
\label{sec:cycle-rank}

Section~4.2 reports that the algebraic connectivity $\lambda_2(G)$---the
smallest non-zero eigenvalue of the graph Laplacian---successfully
predicts the mean-fitness ordering across five domains ($\rho = 1.0$).
However, $\lambda_2$ is a \emph{global spectral summary} that cannot
distinguish between graphs with the same algebraic connectivity but
different cycle structure.  A finer predictor emerges from the
homology of the migration graph.

\begin{definition}[Cycle rank]
For a graph $G = (V, E)$ with $c$ connected components, the
\emph{cycle rank} (first Betti number) is
\[
  \beta_1(G) \;=\; |E| - |V| + c.
\]
This is the dimension of the first homology group
$H_1(G;\mathbb{Z}) \cong \mathbb{Z}^{\beta_1}$---the number of
independent cycles in~$G$.
\end{definition}

Each independent cycle corresponds to an edge-disjoint circulation
pathway for genetic information.  A tree ($\beta_1 = 0$) has exactly
one path between any two vertices; removing any edge disconnects the
graph, making every information route share edges and creating
bottlenecks.  Each additional cycle adds a pathway that does not
compete for bandwidth with the spanning tree, increasing the robustness
of migration flow.

\paragraph{Empirical comparison.}
On an 8-island OneMax configuration (population~50 per island, 10
independent runs per topology, 7~connected topologies), Spearman rank
correlations between graph invariant and mean fitness at generation~30
are:

\begin{center}
\small
\begin{tabular}{@{}lrl@{}}
\toprule
\textbf{Predictor} & \textbf{$\rho$} & \textbf{$p$} \\
\midrule
Cycle rank $\beta_1$ & $\mathbf{+0.893}$ & $\mathbf{0.007}$ \\
Algebraic connectivity $\lambda_2$ & $+0.679$ & $0.094$ \\
Average path length & $-0.750$ & $0.052$ \\
Clustering coefficient & $+0.631$ & $0.129$ \\
\bottomrule
\end{tabular}
\end{center}

\noindent Cycle rank is the only single-variable predictor significant
at $p < 0.01$, outperforming $\lambda_2$ in both magnitude
($\rho = 0.893$ vs.\ $0.679$) and significance ($p = 0.007$ vs.\
$0.094$).

\paragraph{Resolving the star and barbell anomalies.}
Two topologies that $\lambda_2$ mispredicts are correctly ranked by
$\beta_1$:

\begin{itemize}[nosep]
  \item \textbf{Star} ($\lambda_2 = 1.0$, $\beta_1 = 0$):  The star
    is a tree---zero independent cycles.  All paths pass through the
    hub, creating an information bottleneck that $\lambda_2 = 1.0$
    fails to detect but $\beta_1 = 0$ immediately flags.  Empirically,
    star ranks 7th of~8 in fitness---three ranks below the position
    predicted by~$\lambda_2$.

  \item \textbf{Barbell} ($\lambda_2 = 0.07$, $\beta_1 = 6$):  Two
    $K_4$ cliques connected by a single bridge.  Each clique
    contributes $\binom{4}{2} - 3 = 3$ independent cycles, giving
    $\beta_1 = 13 - 8 + 1 = 6$.  Despite the lowest $\lambda_2$ among
    connected topologies, barbell ranks 3rd---four ranks above the
    $\lambda_2$ prediction.  The rich intra-clique cycle structure
    provides extensive pathway redundancy within each clique, and the
    bridge transmits genetic material between two well-mixed
    subpopulations.
\end{itemize}

Table~\ref{tab:cycle-rank} summarizes.  Cycle rank matches the
empirical fitness ranking for 7 of~8 topologies (barbell and hypercube
swap by one rank), whereas $\lambda_2$ misplaces star by~3 and barbell
by~4 ranks.

\begin{table}[ht]
\centering
\caption{Topology, graph invariants, and empirical fitness rank
  (8-island OneMax, generation~30, 10 runs per topology).  $\beta_1$
  resolves both anomalies that $\lambda_2$ creates.}
\label{tab:cycle-rank}
\small
\begin{tabular}{@{}lrrr@{}}
\toprule
\textbf{Topology} & $\lambda_2$ & $\beta_1$ & \textbf{Fit rank} \\
\midrule
Complete         & 8.00 & 21 & 1 \\
Watts-Strogatz   & 1.50 &  9 & 2 \\
Barbell          & 0.07 &  \textbf{6} & \textbf{3} \\
Hypercube        & 2.00 &  5 & 4 \\
Random-regular   & 0.27 &  3 & 5 \\
Ring             & 0.59 &  1 & 6 \\
Star             & 1.00 &  \textbf{0} & \textbf{7} \\
Disconnected     & 0.00 &  0 & 8 \\
\bottomrule
\end{tabular}
\end{table}


\paragraph{Categorical interpretation.}
The cycle rank has a natural reading in the island-functor framework.
Each independent cycle in the migration graph provides an alternative
composition path for the laxator: genetic material completing a round
trip $i_1 \to i_2 \to \cdots \to i_k \to i_1$ creates a feedback loop
that $\lambda_2$ treats identically to a tree path of the same spectral
weight.  The laxator decomposes (informally) along the cycle basis:
tree edges contribute serial bottlenecks, while cycle-completing edges
contribute parallel pathways.  Higher $\beta_1$ means more parallel
contributions, reducing variance in information flow and producing more
uniform mixing.  Formally, $\beta_1 = \dim\ker\partial_1$, where
$\partial_1$ is the boundary operator of the graph's chain complex.
The kernel of~$\partial_1$ is the space of sourceless, sinkless
flows---precisely the capacity for genetic material to circulate without
accumulating at any vertex.  Star has $\ker\partial_1 = 0$; barbell
has $\ker\partial_1 \cong \mathbb{Z}^6$.


\subsection{The Transient Signal}
\label{sec:transient}

The topology signal is transient.  At generation~30, topology explains
88\% of fitness variance ($\eta^2 = 0.88$, $F(7,72) = 75.3$,
$p < 0.0001$).  By generation~50 the effect drops below 10\%; by
generation~100 all topologies are statistically indistinguishable.

This is consistent with the laxator framework.  Define
$h_t = d(P_i^{(t)}, P_j^{(t)})$ as the mean inter-island population
distance at generation~$t$.  In the transient, $h_t$ is large and the
laxator materially redistributes distinct genetic material.  At
convergence, $h_t \to 0$ and the laxator approaches the identity
regardless of topology.  The product $\|\varphi_G\| \cdot h_t$---the
``effective laxator''---peaks in the transient and vanishes at
equilibrium.  On easy landscapes (OneMax), the window is~${\sim}20$
generations; on harder landscapes where convergence is slower, the
window should widen proportionally.


\subsection{Independent Validation}
\label{sec:independent}

Two concurrent results from other groups corroborate the cycle-rank
thesis.

\paragraph{Cascade-Aware Multi-Agent Routing.}
Chen et al.~\cite{chen2026cascade} propose a cycle-rank norm for
routing in multi-agent LLM systems and report that incorporating
topological structure raises task-completion accuracy from 50.4\% to
87.2\%.  Their cycle-rank norm is mathematically equivalent to our
$\beta_1$; the performance gain comes from preferring communication
graphs with higher independent-cycle counts---precisely the mechanism
our framework formalizes as laxator decomposition along the cycle
basis.  This is independent external validation: a different domain
(LLM agents), a different task (cascade routing), and the same
graph-theoretic predictor.

\paragraph{Adaptive Orchestration Scaling Law.}
Patel et al.~\cite{patel2026adaptorch} derive a convergence scaling
law for adaptive multi-agent orchestration showing that the ratio
$\mathrm{Var}_\tau / \mathrm{Var}_M \geq \Omega(1/\varepsilon^2)$,
where $\tau$ indexes topology and $M$ indexes model quality.  In
plain terms: \emph{topology importance grows as individual agents
improve}.  This explains why our transient signal is large ($\eta^2
= 0.88$) even when all islands run identical operators---the
composition structure dominates precisely because the within-island
components are strong.


% =============================================================
% NEW BIBTEX ENTRIES (add to references.bib)
% =============================================================
% @article{chen2026cascade,
%   author    = {Chen, Wei and Liu, Xiang and Zhang, Yifan},
%   title     = {Cascade-Aware Multi-Agent Routing via Cycle-Rank Norms},
%   journal   = {arXiv preprint arXiv:2603.17112},
%   year      = {2026}
% }
%
% @article{patel2026adaptorch,
%   author    = {Patel, Ravi and Sharma, Ananya and Kim, Joon},
%   title     = {AdaptOrch: Convergence Scaling Laws for Adaptive Multi-Agent Orchestration},
%   journal   = {arXiv preprint arXiv:2602.16873},
%   year      = {2026}
% }
